\setcounter{subsection}{0}
\section{DISEÑO DE LA SOLUCIÓN}

Este capítulo presenta el resultado de la Fase~2: la arquitectura de Tukit y las decisiones del
diseño detallado. Su
argumento central es que las fronteras entre servicios se justifican por las propiedades que
los datos no pueden compartir, no por la carga que el sistema va a soportar.

La exposición sigue la norma ISO/IEC/IEEE 42010 \citep{iso_42010_2022} en la separación
entre preocupaciones de los interesados, puntos de vista y vistas, y el modelo 4+1 de
Kruchten \citep{kruchten_41_1995} para distinguir la vista de desarrollo de la vista
física, distinción que resulta central en la Sección~\ref{sec:monorepo}. Las alternativas se
evaluaron con ATAM y las decisiones quedaron en registros de decisión de arquitectura.

% =====================================================================
\subsection{Impulsores arquitectónicos}
\label{sec:impulsores}

El diseño inicial justificaba la descomposición en microservicios por
\emph{escalabilidad}. Un cálculo de carga con las cifras del mercado objetivo (diecisiete
colegios en el corpus y una demanda concentrada en las semanas previas al inicio del año
lectivo) mostró que la infraestructura mínima excede la demanda proyectada en unos tres
órdenes de magnitud. Esa justificación no se sostenía. El rediseño condujo a tres criterios,
ninguno de rendimiento, que resume la Tabla~\ref{tab:impulsores}.

\begin{table}[htbp]
\centering
\caption{Impulsores arquitectónicos de la descomposición en servicios}
\label{tab:impulsores}
\small
\begin{tabular}{@{}p{3.1cm}p{5.0cm}p{5.6cm}@{}}
\toprule
\textbf{Criterio} & \textbf{Fundamento} & \textbf{Qué separa} \\
\midrule
Recuperabilidad asimétrica &
El objetivo de punto de recuperación difiere en órdenes de magnitud entre almacenes &
El catálogo se reconstruye desde sus archivos de datos, una vez exportada al repositorio la
curaduría del panel, y admite perder un día de cambios (CF-02a); un pedido no puede perder un
solo byte. Compartir almacén obliga a dimensionar el
respaldo por el más exigente y a exponer el crítico a la carga de escritura del volátil \\
\addlinespace
Frontera de cumplimiento &
Ley 1581 de 2012 \citep{congreso_ley1581_2012} y Decreto 1377 de 2013
\citep{presidencia_decreto1377_2013} &
Los datos de menores de edad se aíslan en un servicio con base, proceso y llaves propios.
La superficie que un auditor debe examinar se reduce a un solo componente \\
\addlinespace
Perfil de recurso &
Heterogeneidad de la demanda computacional &
La inferencia de modelos y la optimización multiobjetivo son intensivas en CPU y memoria;
la gestión transaccional requiere disponibilidad y consistencia fuerte. Colocarlas juntas
obliga a dimensionar todo por la más exigente \\
\bottomrule
\end{tabular}
\end{table}

Con estos criterios, el número de servicios que ya corría en los prototipos resultó
aproximadamente adecuado; lo que estaba mal era la ubicación de algunas fronteras.

% =====================================================================
\subsection{De catorce servicios documentados a seis desplegables}
\label{sec:consolidacion}

La arquitectura de los documentos académicos iniciales proponía trece microservicios de
dominio y una puerta de enlace, para separar al máximo las responsabilidades. Dos hechos la
cuestionaron. El primero fue de costo y operabilidad: catorce desplegables implican catorce
canales de integración y despliegue, catorce bases candidatas y una malla de red que un
equipo de cuatro personas no puede operar, frente a una demanda proyectada de unos treinta y
seis pedidos diarios en temporada. El segundo fue empírico: la base de grafos Neo4j,
prevista para la sustituibilidad, estaba apagada por configuración y el grafo de conceptos
cabía en memoria, por lo que el ADR 0006 la retiró y resolvió el grafo en proceso; su
evaluación como GraphRAG consta en el anexo CRISP-DM. El sistema construido sustituyó a su
vez ese grafo por una tabla curada y dirigida de sustituciones en el servicio de catálogo,
que tiene modelo, migración y pruebas. Ningún servicio la consulta en la versión evaluada, de modo que la sustitución se delimitó como línea de la evolución del catálogo (Sección~\ref{sec:trabajo-futuro}).

El ADR 0001 aplicó los tres criterios de la Tabla~\ref{tab:impulsores} y redujo los catorce
servicios a seis más la puerta de enlace. Los microservicios anteriores, desarrollados entre
el 2 y el 18 de septiembre de 2026, hicieron esa fusión el 10 de septiembre. Identidad y
privacidad quedaron en un servicio de cuentas; catálogo, pedidos y última milla, en uno de
comercio, y PQRS, analítica y notificaciones, en uno de plataforma. Aparte quedaron
instituciones, una fachada para el panel de administración y el servicio de búsqueda.

Esa reagrupación
dejaba un servicio de respaldo para el panel de administración y un servicio de
instituciones separado de las cuentas. Con la razón de cambio como criterio, los colegios
pasaron a identidad, el catálogo se separó de comercio (ADR 0011), el panel consume la misma
API que la aplicación del acudiente y la auditoría pasó a ser una biblioteca compartida
(ADR 0002). El servicio de búsqueda anterior se dividió entre catálogo e inteligencia,
porque sus objetivos de recuperación no permiten compartir almacén. El resultado es mejor
para un equipo de cuatro personas: hay menos procesos que operar, cada servicio tiene un
solo motivo para cambiar y ningún contrato se duplica. El sistema construido tiene seis
desplegables con código, cinco servicios de dominio (identidad, catálogo, inteligencia,
comercio y operación) y la puerta de enlace, sobre cinco bases.

El diseño de la
Fase 2 preveía un séptimo desplegable, ingesta, que extraía de forma masiva los productos y
precios de las tiendas. El sistema construido lo reemplazó por un catálogo curado que crece
con la demanda: un producto entra cuando una lista lo pide y el equipo confirma la
correspondencia. Es mejor porque concentra el esfuerzo en los productos que las listas
piden, reduce los candidatos de cada concepto de miles a decenas y evita publicar un
producto parecido pero equivocado. En el sistema construido reúne 1.050 productos (1.040
activos), 1.057 ofertas de ocho tiendas y 149 conceptos en archivos de datos versionados, y las herramientas que
consultan las tiendas solo completan o enriquecen productos. Lo que el equipo cura en el
panel queda en la base del catálogo hasta que se exporta a esos archivos (ADR 0028).

% =====================================================================
\subsection{Alternativas evaluadas}
\label{sec:alternativas}

Se compararon tres alternativas mediante el método ATAM \citep{bass_software_2021}, cuyos
resultados completos se consignan en el documento de evaluación ATAM. La
Tabla~\ref{tab:alternativas} sintetiza el contraste.

\begin{table}[htbp]
\centering
\caption{Alternativas arquitectónicas evaluadas}
\label{tab:alternativas}
\small
\begin{tabular}{@{}p{2.6cm}p{3.5cm}p{3.5cm}p{4.1cm}@{}}
\toprule
& \textbf{A. Monolito modular} & \textbf{B. Microservicios de grano fino} & \textbf{B$'$. Grano grueso (elegida)} \\
\midrule
Componentes & 1 desplegable, módulos internos & 13 servicios y puerta de enlace & 5 servicios de dominio y puerta de enlace, 5 bases \\
\addlinespace
Aislamiento de datos de menores & Por convención de código & Por servicio & Por servicio, base y llaves \\
\addlinespace
Punto de recuperación & Único para todo & Por servicio & Diferenciado por criticidad \\
\addlinespace
Complejidad operativa & Baja & Alta: coordinación distribuida en cada transacción & Media \\
\addlinespace
Costo de infraestructura & Mínimo & Alto & Moderado \\
\addlinespace
Puntaje ponderado (sobre 5,00) & 2,82 & 3,44 & 3,91 \\
\bottomrule
\end{tabular}
\end{table}

La alternativa~A se descarta por el primer criterio de la Tabla~\ref{tab:impulsores}: un
almacén único impone un mismo objetivo de recuperación a datos cuyo valor difiere en
órdenes de magnitud. La alternativa~B se descarta por complejidad operativa: fragmentar la
gestión de pedidos, inventario y pago en servicios separados convierte cada compra en una
transacción distribuida, y el equipo no dispone de los recursos de operación que eso exige
\citep{newman_building_2021}. La alternativa~B$'$ conserva el aislamiento donde el
cumplimiento y la recuperabilidad lo exigen, y mantiene junta la transaccionalidad que
debe permanecer junta. La matriz ponderada del ATAM le da 3,91 sobre 5,00, frente a 3,44 de
la granularidad fina y 2,82 del monolito modular.

La Figura~\ref{fig:c4-contexto} (Anexo~\ref{anx:figuras}) sitúa el sistema frente a sus actores y sistemas externos, y
la Figura~\ref{fig:c4-contenedores} materializa la alternativa elegida: cinco servicios de
dominio y la puerta de enlace sobre cinco bases de datos, con una bandeja de salida en cada
base para entregar los eventos. Las dos salen del código fuente: la primera es una vista UML
del modelo de arquitectura, y la segunda sigue el modelo C4. La vista de contenedores añade los procesos en segundo plano de cada servicio y numera
los pasos de tres flujos: la subida de la lista y los kits, el pago y los avisos.

Con RabbitMQ, un publicador por servicio productor lleva los eventos de su bandeja a una cola
por consumidor, con una cola de reintentos y otra de descartados, y la entrega es al menos una
vez, con consumidores idempotentes (ADR 0037). La bandeja sigue siendo la fuente de verdad. Sin
RabbitMQ, cada consumidor lee la bandeja en la base del productor, sin tope de intentos ni cola
de descartados. El manifiesto declara esas lecturas como respaldo, y una función de aptitud
impide que el camino con RabbitMQ lea bases ajenas.



\begin{figure}[htbp]
\centering
% Vista de contenedores (C4 nivel 2) del código fuente del sistema en su versión del 24 de octubre de 2026,
% con los flujos de la subida y los kits (A), el pago (B) y los avisos (C). Estilos locales: no depende del
% preámbulo; solo usa las bibliotecas arrows.meta y shapes.geometric.
\begin{tikzpicture}[x=1cm, y=1cm, font=\scriptsize,
  per/.style={draw, rounded corners=7pt, align=center, fill=yellow!10, inner sep=3pt, minimum height=7mm},
  cont/.style={draw, rounded corners=2pt, align=center, fill=blue!8, inner sep=2pt, minimum height=8mm},
  proc/.style={draw, rounded corners=2pt, align=center, fill=green!10, inner sep=2pt, minimum height=7mm},
  ext/.style={draw, rounded corners=2pt, align=center, fill=orange!14, inner sep=2pt, minimum height=7mm},
  infra/.style={draw, rounded corners=2pt, align=center, fill=gray!12, inner sep=2pt, minimum height=7mm},
  bd/.style={draw, cylinder, shape border rotate=90, aspect=0.22, align=center, fill=white,
    minimum width=1.25cm, minimum height=8mm, inner sep=1pt},
  grupo/.style={draw=gray!70, rounded corners=3pt, fill=gray!3},
  http/.style={-{Stealth[length=1.6mm]}, semithick},
  evt/.style={-{Stealth[length=1.6mm]}, densely dashed, semithick},
  cola/.style={-{Stealth[length=1.6mm]}, semithick, draw=violet!75!black},
  hueco/.style={draw=white, line width=2.8pt},
  paso/.style={circle, draw, fill=white, inner sep=0.5pt, font=\tiny\bfseries, minimum size=3.6mm},
  rot/.style={font=\tiny, fill=white, inner sep=1pt},
  col/.style={anchor=north west, text width=4.9cm, align=left, inner sep=0pt, font=\scriptsize}]
% Personas y aplicaciones
\node[per, text width=2.6cm] (acudiente) at (2.45,0) {\textbf{Acudiente}};
\node[per, text width=3.8cm] (equipo) at (8.6,0) {\textbf{Equipo de Tukit}\\operador y administrador};
\node[cont, text width=3.4cm] (web) at (2.45,-1.45) {\textbf{Aplicación web progresiva}\\Next.js, puerto 3000};
\node[ext, text width=1.85cm] (google_places) at (5.5,-1.45) {\textbf{Google Maps}\\sugerencias de dirección};
\node[cont, text width=3.4cm] (admin) at (8.6,-1.45) {\textbf{Panel de administración}\\Next.js, puerto 3001};
\node[ext, text width=2.5cm] (wompi) at (13.4,-1.45) {\textbf{Wompi}\\pasarela de pago};
\node[cont, minimum width=14.3cm, text width=14cm] (gateway) at (7.35,-2.85) {\textbf{Puerta de enlace}
  (Python, puerto 8010): reenvía cada petición al servicio que atiende el prefijo de su ruta};
% Servicios que llaman a otros: operación, inteligencia y comercio
\draw[grupo] (0.2,-3.65) rectangle (4.55,-6.25);
\draw[grupo] (4.85,-3.65) rectangle (10.15,-6.25);
\draw[grupo] (10.45,-3.65) rectangle (14.95,-6.25);
\node[cont, text width=2.3cm] (operacion) at (1.6,-4.3) {\textbf{Operación}\\puerto 8015};
\node[bd] (base_operacion) at (3.7,-4.3) {base de\\operación};
\node[proc, text width=1.7cm] (instantanea) at (1.2,-5.55) {instantánea de métricas (a mano)};
\node[proc, text width=1.7cm] (repartidor) at (3.5,-5.55) {repartidor de avisos};
\node[cont, text width=2.1cm] (inteligencia) at (6.4,-4.3) {\textbf{Inteligencia}\\puerto 8013};
\node[bd] (base_inteligencia) at (9.2,-4.3) {base y cola\\de trabajos};
\node[proc, text width=1.45cm] (sondeador_listas) at (5.75,-5.55) {sondeador\\de listas};
\node[proc, text width=1.45cm] (trabajador) at (7.5,-5.55) {trabajador\\durable};
\node[proc, text width=1.45cm] (publicador_inteligencia) at (9.25,-5.55) {publicador\\de eventos};
\node[cont, text width=2.4cm] (comercio) at (11.85,-4.3) {\textbf{Comercio}\\puerto 8014};
\node[bd] (base_comercio) at (13.95,-4.3) {base de\\comercio};
\node[proc, text width=1.6cm] (barredor) at (11.5,-5.55) {barredor de reservas};
\node[proc, text width=1.45cm] (publicador_comercio) at (13.75,-5.55) {publicador\\de eventos};
% Bus de llamadas a identidad y catálogo, que no llaman a otros servicios
\draw[semithick] (2.95,-6.7) -- (12.55,-6.7);
\foreach \x in {2.95,8.55,12.55}{\draw[semithick] (\x,-6.25) -- (\x,-6.7);}
\draw[semithick] (10.3,-3.25) -- (10.3,-6.7);
\node[rot, align=center] at (11.42,-6.7) {llamadas HTTP a\\identidad y catálogo};
% RabbitMQ: intercambio de eventos entre las dos filas de servicios
\node[draw=violet!75!black, fill=violet!10, rounded corners=2pt, minimum width=10.35cm, minimum height=6.5mm,
  inner sep=0pt] (rabbitmq) at (9.125,-7.25) {};
\node[align=center, inner sep=0pt, font=\scriptsize] at (5.95,-7.25) {\textbf{RabbitMQ}\\[-1.5pt]intercambio tukit.eventos};
% Servicios proveedores
\draw[grupo] (2.9,-7.8) rectangle (6.75,-10.3);
\node[cont, text width=1.8cm] (identidad) at (3.95,-8.45) {\textbf{Identidad}\\puerto 8011};
\node[bd] (base_identidad) at (6.0,-8.45) {base de\\identidad};
\node[proc, text width=2.4cm] (sondeador_colegios) at (4.75,-9.65) {sondeador de colegios};
\draw[grupo] (9.05,-7.8) rectangle (13.1,-10.3);
\node[cont, text width=2.0cm] (catalogo) at (10.2,-8.45) {\textbf{Catálogo}\\puerto 8012};
\node[bd] (base_catalogo) at (12.3,-8.45) {base de\\catálogo};
\node[proc, text width=2.0cm, fill=green!4, densely dashed] (herramientas) at (10.2,-9.65) {herramientas de\\curaduría (a mano)};
\node[proc, text width=1.3cm] (publicador_catalogo) at (12.15,-9.65) {publicador\\de eventos};
% Sistemas externos
\node[ext, text width=1.9cm] (gemini) at (7.85,-8.45) {\textbf{Gemini}\\lectura y vectores};
\node[ext, text width=1.7cm] (tiendas) at (14.4,-9) {\textbf{Tiendas en línea}\\productos y precios};
\node[ext, text width=1.9cm] (correo) at (1.25,-8.45) {\textbf{Correo saliente}\\SMTP};
\node[ext, text width=1.9cm] (whatsapp) at (1.25,-9.55) {\textbf{WhatsApp}\\enlace wa.me};
% Fila de infraestructura
\node[infra, text width=4.6cm] (postgres) at (4.9,-11.05) {\textbf{PostgreSQL 16 con pgvector}\\un servidor con las
  cinco bases;\\cada una con su bandeja de salida};
\node[infra, text width=4.2cm, densely dotted, anchor=north] (redis) at (9.6,-10.45) {\textbf{Redis}, opcional\\caché de vectores y de lecturas\\y contador del límite de tasa\\de la puerta de enlace};
\node[infra, text width=3.3cm, densely dotted] (observabilidad) at (13.55,-11.05) {\textbf{Observabilidad opcional}\\OpenTelemetry y Jaeger};
% Personas, aplicaciones y pasarela
\draw[http] (acudiente) -- (web);
\draw[http] (equipo) -- (admin);
\draw[http] (web) -- (google_places);
\draw[http] (web.south) -- (web.south |- gateway.north);
\draw[http] (admin.south) -- (admin.south |- gateway.north);
\draw[http] ([xshift=2mm]web.north west) -- (0.88,0.65) -| node[rot, pos=0.25, above] {abre el widget de pago de Wompi} (wompi.north);
\draw[http] (wompi.south) -- (wompi.south |- gateway.north) node[rot, midway, right] {aviso firmado};
\draw[http] ([xshift=-2mm]comercio.north east) -- ++(0,0.45) -- (15.25,-3.45) |- (wompi.east);
\draw[http] (gateway.south -| operacion) -- (operacion.north);
\draw[http] (gateway.south -| inteligencia) -- (inteligencia.north);
\draw[http] (gateway.south -| comercio) -- (comercio.north);
% Llamadas entre servicios que llaman
\draw[http] ([xshift=-2mm]operacion.north east) -- ++(0,0.35) -| ([xshift=3mm]inteligencia.north west);
\draw[http] ([xshift=-3mm]inteligencia.north east) -- ++(0,0.4) -| ([xshift=3mm]comercio.north west);
% Llamadas a identidad y catálogo desde el bus; la del catálogo cruza RabbitMQ
\draw[http] (3.6,-6.7) -- (3.6,0 |- identidad.north);
\draw[hueco] (10.3,-6.9) -- (10.3,-7.6);
\draw[http] (10.3,-6.7) -- (10.3,0 |- catalogo.north);
% Bandejas de salida que lee cada publicador
\draw[evt] (base_inteligencia.south) -- (base_inteligencia.south |- publicador_inteligencia.north);
\draw[evt] (base_comercio.south) -- (base_comercio.south |- publicador_comercio.north);
\draw[evt] (base_catalogo.south) -- (base_catalogo.south |- publicador_catalogo.north);
% Eventos por RabbitMQ: de los publicadores al intercambio y del intercambio a la cola de cada consumidor
\foreach \x in {4.1,5.75,7.1,9.25}{\draw[hueco] (\x,-6.58) -- (\x,-6.82);}
\draw[cola] (publicador_inteligencia.south) -- (9.25,0 |- rabbitmq.north);
\draw[cola] (publicador_comercio.south) -- (13.75,0 |- rabbitmq.north);
\draw[cola] (11.5,0 |- publicador_catalogo.north) -- (11.5,0 |- rabbitmq.south);
\draw[cola] (4.1,0 |- rabbitmq.north) -- (4.1,0 |- repartidor.south);
\draw[cola] (5.75,0 |- rabbitmq.north) -- (sondeador_listas.south);
\draw[cola] (7.1,0 |- rabbitmq.north) -- (7.1,0 |- trabajador.south);
\draw[cola] (5.2,0 |- rabbitmq.south) -- (5.2,0 |- sondeador_colegios.north);
% Llamadas a sistemas externos
\draw[hueco] (7.85,-6.58) -- (7.85,-6.82);
\draw[hueco] (7.85,-6.9) -- (7.85,-7.6);
\draw[http] (7.85,0 |- trabajador.south) -- (gemini.north);
\draw[http] (catalogo.west) -- (gemini.east);
\draw[http] (13.1,-9) -- (tiendas.west);
\draw[http] (2.75,0 |- repartidor.south) |- (correo.east);
\draw[semithick] (2.75,0 |- repartidor.south) |- (whatsapp.east);
% Pasos de los flujos
\node[paso] at (2.15,-2.15) {A1};
\node[paso] at (7.85,-7.73) {A2};
\node[paso] at (8.3,-2.15) {A3};
\node[paso] at (5.75,-6.45) {A4};
\node[paso] at (2.75,-2.15) {B1};
\node[paso] at (11.6,0.65) {B2};
\node[paso] at (13.05,-2.15) {B3};
\node[paso] at (barredor.north west) {B4};
\node[paso] at (4.1,-6.45) {C1};
\node[paso] at (2.95,-6.45) {C2};
\node[paso] at (2.75,-9) {C3};
% Pasos de los flujos y notación, debajo de la figura
\node[col] (colA) at (0.2,-12.1) {\textbf{A. Subida de la lista y kits}\\[1pt]
  \textbf{A1} El acudiente sube la lista; inteligencia guarda el trabajo y lo anuncia en su bandeja de salida.\\[1pt]
  \textbf{A2} El trabajador toma el trabajo de RabbitMQ o de su tabla, lee las páginas con Gemini, pide candidatos al catálogo y deja los kits por revisar.\\[1pt]
  \textbf{A3} El equipo revisa los kits en el panel y los publica; inteligencia deja el evento de kits publicados en su bandeja de salida.\\[1pt]
  \textbf{A4} Si la lista no estaba verificada, el sondeador de listas arma el kit al recibir el evento de lista verificada.};
\node[col] (colB) at (5.35,-12.1) {\textbf{B. Pago}\\[1pt]
  \textbf{B1} El acudiente elige un kit y escribe la dirección con sugerencias de Google Maps; inteligencia crea la orden en comercio, que reserva el inventario al precio vigente.\\[1pt]
  \textbf{B2} La aplicación web abre el widget de Wompi y el acudiente paga.\\[1pt]
  \textbf{B3} Wompi envía el aviso firmado o la página de resultado pide sincronizar el pago; comercio deja el evento de pedido pagado en su bandeja.\\[1pt]
  \textbf{B4} El barredor libera la reserva del pedido que no se paga en 24 horas y lo cancela.};
\node[col] (colC) at (10.5,-12.1) {\textbf{C. Avisos}\\[1pt]
  \textbf{C1} El repartidor de avisos recibe por RabbitMQ los eventos de kits publicados, pago, incidencia, despacho, reparto y entrega, y no repite el aviso de un evento ya enviado.\\[1pt]
  \textbf{C2} Pide a identidad el contacto del acudiente, sin guardarlo.\\[1pt]
  \textbf{C3} Si el acudiente tiene celular, arma un enlace de WhatsApp con el aviso; si no, se lo envía por correo.\\[1pt]
  El despacho lo registra el equipo en el panel, con el nombre y el celular del repartidor propio; las transportadoras con guía son una evolución prevista.};
\node[anchor=north west, text width=15.1cm, align=left, inner sep=0pt, font=\scriptsize]
  at ([yshift=-2mm]colA.south west |- current bounding box.south) {Línea continua: llamada HTTP. Línea discontinua:
  el publicador lee la bandeja de salida de su base. Línea violeta: evento que viaja por RabbitMQ hasta la cola de su
  consumidor (avisos, colegios, listas y trabajos), con reintentos y descartados; sin RabbitMQ, los publicadores no
  arrancan y cada consumidor lee la bandeja del productor. Línea punteada: componente opcional. Línea sin punta: enlace que el sistema
  arma y no envía. Recuadro gris: el servicio con sus procesos y su base.};
\end{tikzpicture}

\caption{Vista de contenedores (C4 nivel 2) del sistema construido, con los procesos en
segundo plano de cada servicio y los pasos de la subida de la lista y los kits (A), el pago (B)
y los avisos (C).}
\label{fig:c4-contenedores}
\end{figure}

Redis es opcional y solo guarda lo que se puede recalcular o perder (ADR 0038): el vector de cada
renglón durante 24 horas, la lectura aceptada de cada lista durante 30 días y los contadores del
límite de tasa de la puerta de enlace. Si no responde en 250~ms, cada servicio calcula sin caché
y la puerta de enlace cuenta el límite en su memoria.

% =====================================================================
\subsection{Modelo de dominio y contextos delimitados}
\label{sec:dominio}

La descomposición sigue el enfoque de diseño guiado por el dominio \citep{evans_domain-driven_2004}.
Cada servicio corresponde a un contexto delimitado con un lenguaje ubicuo propio, y las
traducciones entre contextos ocurren explícitamente en los adaptadores.

Las migraciones del sistema construido crean 74 tablas en cinco bases de datos, una por
servicio (Tabla~\ref{tab:reparto}); la bandeja de salida de eventos se repite en las cinco,
de modo que hay 70 nombres distintos. El recuento cambió durante la implementación, entre
67 y 74 tablas según la versión.

% Tabla tab:reparto movida al Anexo A.3 (tablas del cuerpo)


\subsubsection{La matrícula como ancla de tenencia}

El modelo de datos de los microservicios anteriores no sabía a qué colegio pertenecían los
pedidos ni los usuarios: el identificador de institución aparecía cero veces en los
servicios de comercio y de cuentas. Donde sí existía, lo enviaba el cliente en el cuerpo de
la petición sin validación alguna, lo cual es peor que su ausencia porque aparenta un
control que no existe.

La solución aparentemente obvia, un identificador de colegio en el usuario, es incorrecta,
porque un acudiente puede tener hijos en colegios distintos (CO-04). La tenencia reside, en
cambio, en la relación $\langle$estudiante, sede, grado, año lectivo$\rangle$, que se modela
como la entidad matrícula, con una relación de muchos a muchos con los acudientes
(ADR 0012). Su identificador viaja a los otros cuatro servicios como referencia lógica y es
el eje de las políticas de visibilidad. Se ubica en la \emph{línea} del pedido y no en su
cabecera, porque en la cabecera dos hermanos obligarían a dos pedidos y a dos pagos, lo que
contradice la propuesta de valor del producto. La comprobación de que cada línea apunta a
una matrícula visible para el acudiente se delimitó como parte del endurecimiento productivo
(Sección~\ref{sec:trabajo-futuro}).

\subsubsection{Asimetría entre lectura y escritura}

Las políticas de visibilidad son deliberadamente asimétricas. Para \emph{leer} un recurso
basta que toque alguna matrícula visible para el solicitante; para \emph{escribirlo} se
exige que \emph{todas} las matrículas que el recurso referencia estén contenidas en su
ámbito. Un pedido con dos hermanos, uno visible y otro no, puede consultarse pero no
modificarse. La simetría ingenua produciría, en ese escenario, una fuga de información o un
sistema inutilizable.

% =====================================================================
\subsection{Autorización con ámbito y seguridad a nivel de campo}
\label{sec:autorizacion}

El modelo de autorización de los microservicios anteriores representaba el rol como una
cadena de texto en la tabla de usuarios, y el control de acceso se reducía a verificar que
esa cadena perteneciera a un conjunto permitido. Dos usuarios con el mismo rol sobre
colegios distintos superaban, así, exactamente el mismo control. La única defensa consistía
en que cada punto de entrada recordara comparar el colegio del token con el del recurso
solicitado, y la auditoría de los 152 puntos de entrada de esos servicios encontró 22 en los
que esa comparación faltaba o era parcial (Sección~\ref{sec:defectos}).

El diagnóstico requiere precisión. Uno de esos servicios, el de instituciones, ya definía
una relación de pertenencia con la tripla \emph{(usuario, institución, rol)}, introducida
para corregir una referencia directa insegura real. Expresar autoridad con ámbito no era,
por tanto, imposible: el defecto era de alcance y de completitud. El mecanismo cubría un
servicio de siete; admitía un único nivel de ámbito, la institución, cuando el recurso
protegido pertenece a una sede y un grado; permitía un solo rol por institución, y carecía
de vigencia temporal. Esta última carencia producía una contradicción ilustrativa: la
documentación del modelo justificaba conservar el vínculo para auditar quién tuvo acceso y
cuándo, pero la única forma de revocarlo era eliminarlo. El diseño generaliza esa solución,
ya validada contra un ataque real, con la jerarquía de ámbitos, la vigencia temporal y la
aplicación uniforme a todos los servicios.

Con el control de acceso basado en roles con ámbito, la pregunta que el sistema responde
deja de ser \emph{qué es este usuario} y pasa a ser \emph{qué puede hacer este usuario sobre
este recurso concreto}:

\begin{equation}
\mathrm{puede}(u, p, a) \iff \exists\, r \in \mathcal{R}_{\text{vig}}(u, a) : (r, p) \in \mathcal{P}
\end{equation}

\noindent donde $\mathcal{P}$ es el conjunto de pares de rol y permiso, y
$\mathcal{R}_{\text{vig}}(u, a)$ el de roles vigentes del usuario $u$ sobre el ámbito $a$ o
sobre cualquier ámbito que lo contenga, según la jerarquía
$\text{global} \succ \text{colegio} \succ \text{sede} \succ \text{matrícula}$. La contención
opera en un solo sentido: un rol otorgado sobre un colegio alcanza sus sedes, pero ser
acudiente de una matrícula no confiere autoridad sobre la sede.

\subsubsection{La función falla cerrada}

Las políticas de seguridad a nivel de fila de los microservicios anteriores iniciaban su
condición base verificando si el contexto de usuario era nulo, de modo que toda conexión sin
contexto establecido veía la tabla completa. El control fallaba \emph{abierto}, y su modo de
fallo más probable, omitir el contexto, era precisamente el que concedía acceso total. La
función aquí especificada no tiene ninguna rama que conceda por ausencia de información.

En el sistema construido, identidad guarda las asignaciones de rol con su vigencia y los
pares de rol y permiso, y evalúa los permisos con ámbito (ADR 0012); las rutas protegidas
de inteligencia fallan cerradas (SG-06). La aplicación uniforme (AD-03 y SG-02) se delimitó como
endurecimiento productivo (Sección~\ref{sec:trabajo-futuro}): solo una ruta de producción consulta la función de permiso, y varias rutas que
devuelven datos personales no exigen sesión ni rol. Tampoco hay políticas de seguridad a
nivel de fila, de modo que la evaluación local del ADR 0015 sigue siendo diseño
(Sección~\ref{sec:defectos}).

\subsubsection{Dos ejes independientes, no uno}

La seguridad a nivel de campo usa el mismo mecanismo, con un componente de campo en el
permiso: leer un estudiante y leer su nombre son permisos distintos. El rediseño inicial
afirmó que esos permisos se derivaban \emph{automáticamente} del metadato de clasificación de
datos, y el análisis del registro de tratamiento mostró que no es así. La clasificación mide
sensibilidad y el permiso expresa quién puede ver qué: son dos ejes independientes. De la
clasificación solo se deriva el valor por omisión, de modo que todo campo clasificado como
dato de menor o como personal sensible nace denegado y todo campo público nace permitido.
Colapsar ambos ejes conduce a implementar la mitad del mecanismo creyendo que está completo.

% =====================================================================
\subsection{Protección de datos personales desde el diseño}
\label{sec:datos}

El sistema trata datos de menores de edad, datos personales de los acudientes, como el
contacto y la dirección de entrega, y datos de las instituciones educativas. La Ley 1581 de
2012 \citep{congreso_ley1581_2012} rige los dos primeros, y los terceros cuando identifican
a personas naturales: exige autorización previa y una finalidad legítima informada al
titular, y reconoce el derecho de supresión.

Se adoptó un esquema de cinco niveles de clasificación (público, interno, personal,
personal sensible y dato de menor), aplicado campo por campo en el registro de tratamiento
(Anexo~M), que cubre las cinco bases del sistema construido. Su primera versión, sobre los
microservicios anteriores, revisó 398 columnas de 50 tablas y encontró 126 campos con datos
personales, 20 de ellos clasificados como dato de menor.

\subsubsection{Políticas que el esquema anterior no podía sostener}

El ejercicio reveló un patrón recurrente: políticas correctamente enunciadas en la
documentación que el esquema hacía inejecutables. La Tabla~\ref{tab:politicas} recoge los
casos y su corrección.

% Tabla tab:politicas movida al Anexo A.3 (tablas del cuerpo)


El patrón tiene valor metodológico más allá de este proyecto: una política de tratamiento
de datos es verificable solo si el modelo lógico expone las columnas que la hacen
computable. Documentar la política sin instrumentar el esquema produce cumplimiento
aparente. El riesgo inverso también existe: el sistema construido tiene esas columnas, pero
la supresión propagada existe en el modelo y no como flujo de extremo a extremo
(Sección~\ref{sec:defectos}). El SRS v6.1 deja la supresión propagada, el panel de privacidad
y la retención automática fuera del alcance de este trabajo, y atiende los derechos del
titular mediante un caso de soporte.

\subsubsection{Minimización por columna: la dirección de entrega}

Los microservicios anteriores no tenían ninguna dirección de entrega, en un negocio cuya
propuesta de valor consiste en entregar el kit en el domicilio. Tampoco figuraba en el
esquema objetivo, de modo que la omisión se arrastraba del código al diseño: un defecto por
omisión sobrevive más que uno por comisión, porque nada lo señala.

El dato es sensible de forma poco habitual: más que información de contacto, es la
ubicación física del hogar de una familia con menores, cruzada con el colegio al que
asisten. El ADR 0014 distingue dos cosas que se parecen sin ser la misma: la libreta de
direcciones reutilizables del acudiente, en identidad, editable y sujeta a supresión, y la
copia congelada en el envío, en comercio, inmutable una vez confirmado el pedido. El
razonamiento es análogo al de congelar el precio en la línea del pedido: un hecho histórico
y un dato editable son entidades distintas aunque se escriban igual. Si el envío
referenciara la libreta, corregir un número de apartamento reescribiría el destino de una
entrega pasada.

La partición de columnas responde al principio de minimización. La ciudad y la localidad
quedan en claro porque de ellas depende la asignación de la zona de reparto, mientras que
la calle, el número y el teléfono se cifran en la aplicación: planificar una ruta requiere
conocer la zona, no el apartamento. En el sistema construido esa partición rige para la
libreta, que guarda cifrados la dirección y el contacto. La copia del envío conserva el
destinatario y la dirección en claro, como prueba ante un reclamo de la Ley 1480 de 2011, de
modo que la minimización por columna no cubre todavía el dato congelado.

% =====================================================================
\subsection{Organización del código y funciones de aptitud}
\label{sec:monorepo}

\subsubsection{Un repositorio no es un desplegable}

La implementación se organiza en un monorrepositorio con seis desplegables independientes
(ADR 0013). La decisión exige justificación porque contradice la asociación habitual entre
microservicio y repositorio, asociación que carece de fundamento normativo. La norma
ISO/IEC/IEEE 42010 \citep{iso_42010_2022} compone una descripción arquitectónica de vistas
que responden a preocupaciones distintas, y el modelo 4+1 \citep{kruchten_41_1995} separa
la \emph{vista de desarrollo}, que organiza módulos y bibliotecas, de la \emph{vista
física}, que asigna el software a nodos de ejecución. Dónde vive el código y cómo se
despliega son decisiones de vistas diferentes; en la terminología del modelo C4, un
contenedor no equivale a un repositorio.

La motivación concreta es medible. En los microservicios anteriores, trasladar una función
de redacción de datos a la biblioteca compartida exigió modificar siete archivos de
configuración de dependencias y publicar una etiqueta de versión antes de que cualquier
servicio pudiera instalarla. En el monorrepositorio es un único cambio.

\subsubsection{Fronteras verificadas, no declaradas}

La objeción legítima al repositorio único es que facilita la erosión de las fronteras. La
respuesta de este diseño es instrumentarlas mediante funciones de aptitud
\citep{ford_building_2017}: pruebas automáticas y continuas sobre una característica
arquitectónica. Un manifiesto declara, para cada servicio, su lenguaje, su base de datos, las tablas de las que es
dueño y las referencias que cruzan hacia otros dominios, y doce pruebas lo verifican en cada
cambio: que ningún servicio importe código de otro ni declare una tabla ajena, que toda referencia
cruzada esté declarada, que ninguna tabla tenga dos dueños y que el grafo de dependencias sea
acíclico. Dos pruebas más, añadidas con las colas, y una con Redis, extienden la verificación a
esos canales. La lista completa de las doce pruebas está en el Anexo~\ref{anx:ampliaciones}.
El motivo de instrumentarlas es que el defecto que encontró el rediseño, la tabla de productos
declarada como maestra en dos servicios a la vez, no se produjo por falta de criterio sino porque
nada lo verificaba de forma continua.

Una función de aptitud que nunca ha fallado no está demostrada, pues podría no estar
midiendo nada. Por ello el conjunto incluye un segundo nivel, con veintiocho pruebas, que
suministra a los detectores infracciones plantadas y verifica que las detecten, además de
casos que no deben marcarse. Entre esas infracciones están una tabla ajena declarada como
el doble maestro y el ciclo de dependencias que existió. Ante una infracción plantada, la
integración falla e indica el archivo, la línea, el dominio invadido y la remediación
sugerida.

% =====================================================================
\subsection{Asignación de lenguajes: una decisión medida}
\label{sec:lenguajes}

La hipótesis de partida sostenía que los componentes intensivos en cómputo debían reescribirse en
un lenguaje compilado. Para contrastarla se perfiló una lista escolar de veinte ítems, la mediana
del corpus real, contra un catálogo de 3.187 productos. Las cifras son reportadas por la
documentación del código, no recalculadas. El 95,3\,\% de la latencia lo producía una sola línea,
que reconstruía un grafo de sustituibilidad en cada ítem; construirlo cuesta 1.673~ms y consultarlo,
0,137~ms. Corregido ese defecto, la reducción que aporta cambiar de lenguaje queda, según la ley de
Amdahl \citep{amdahl_validity_1967}, entre el 20,0\,\% y el 21,7\,\%, y el 98,5\,\% de lo que
resta corresponde al optimizador NSGA-II sobre pymoo \citep{blank_pymoo_2020}, sin equivalente
maduro en un lenguaje compilado. Por eso los seis desplegables están en Python, y el manifiesto
declara el lenguaje de cada servicio para que incorporar otro sea una adición. El cálculo
completo y el desglose de latencia (Tabla~\ref{tab:latencia}) están en el
Anexo~\ref{anx:ampliaciones}.

% =====================================================================
\subsection{Decisiones registradas}
\label{sec:adr}

Las decisiones arquitectónicas se documentan como registros en el formato propuesto por
Nygard \citep{nygard_documenting_2011} (título, contexto, decisión, estado y consecuencias),
al que el proyecto añade un apartado de alternativas consideradas. La entrega reúne
treinta y ocho registros, numerados de 0001 a 0038, en el registro de decisiones de
arquitectura (Tabla~\ref{tab:anexos}); la Tabla~\ref{tab:adr} destaca los que estructuran el
diseño de este capítulo.

El ADR 0015 ilustra una propiedad del método. Nació de una contradicción entre decisiones previas
que solo se hizo visible al redactar el criterio de aceptación de un requisito: la política de
seguridad no podía evaluarse en ninguno de los tres servicios donde hacía falta, porque la
matrícula vive en la base de identidad. Un registro obliga a escribir el fundamento donde puede
contrastarse con otro.

% Tabla tab:adr movida al Anexo A.3 (tablas del cuerpo)


% =====================================================================
\subsection{Síntesis}

El diseño presentado se distingue del original en tres aspectos, que constituyen el aporte
de ingeniería de este objetivo específico. El primero es la justificación de la
descomposición, que pasó de la escalabilidad, refutada por el cálculo de carga, a la
recuperabilidad asimétrica, la frontera de cumplimiento y el perfil de recurso. Con ella,
los catorce servicios documentados se redujeron a seis desplegables con código, cuyas
fronteras se trazaron de nuevo. El segundo es la matrícula, una entidad que el modelo
anterior no tenía y que permite expresar tenencia, autoridad con ámbito y revocación
parcial del consentimiento. El tercero es que las fronteras entre servicios dejaron de ser
afirmaciones y pasaron a ser aserciones verificables de forma continua. La autorización
uniforme y la depuración del nombre del menor en el documento cargado se delimitaron como
endurecimiento productivo (Sección~\ref{sec:trabajo-futuro}).
